\section{Введение}

\begin{frame}{Темы лекции}
  \begin{enumerate}
    \item Задачи и способы обучения.
    \item Метрики качества.
    \item Признаки и проблемы в данных.
    \item Обучение, недообучение и переобучение.
  \end{enumerate}
  \medskip
  Как проверить, что модель поможет на новых данных?
\end{frame}

\begin{frame}{Когда трудно написать правила вручную?}
  Как написать программу, которая отличает кошку от собаки на фотографии?
  \begin{block}{Что усложняет задачу}
    Разные позы, освещение, фон и породы. Список ручных правил быстро становится слишком большим.
  \end{block}
  \medskip
  Назовите ещё две задачи, в которых примеры получить проще, чем описать все правила.
\end{frame}

\begin{frame}{Машинное обучение: правила из примеров}
  Модель учится отличать кошку от собаки по размеченным фотографиям.
  \begin{description}
    \item[Данные] Фотографии и известные ответы.
    \item[Модель] Функция, которая по фотографии предсказывает класс.
    \item[Обучение] Подбор параметров по размеченным примерам.
  \end{description}
  \medskip
  Закономерность в обучающей выборке может не работать на новых данных.
\end{frame}

\begin{frame}{Объект, признаки и ответ}
  Объект --- квартира; признаки $x$ --- площадь, район и этаж.
  Целевая величина $y$ (\emph{таргет}) --- известная цена продажи.
  \[
    x \quad\xrightarrow{\text{модель } f}\quad \widehat y=f(x)
  \]
  \begin{itemize}
    \item $y$ --- известный ответ; $\widehat y$ --- предсказание.
    \item Примеры $(x_i,y_i)$ образуют размеченную выборку, или датасет.
    \item Для обучения нужны признаки и ответы; для предсказания --- признаки нового объекта.
  \end{itemize}
\end{frame}
